\chapter{Valutazione sperimentale e discussione}
\label{chap:experiments}

% STRUTTURA PROVVISORIA: definire protocollo, baseline, ablation, verifiche
% di generalizzazione, metriche e logging dopo la stesura dei capitoli 1--3.
% Le sezioni indicano domande da affrontare, non risultati gia disponibili.

\section{Obiettivi della valutazione e protocollo sperimentale}
\label{sec:esperimenti-protocollo}
% Collegamento alla domanda di ricerca, scenari, condizioni di confronto,
% metriche, ripetizioni e trattamento dei fallimenti.

\section{Caratterizzazione del costo energetico sui diversi terreni}
\label{sec:esperimenti-costo-terreni}
% Differenze di costo e dipendenza dalla velocita per motivare e
% interpretare gli scenari di scelta del percorso.

\section{Analisi della rappresentazione visiva dei terreni}
\label{sec:esperimenti-feature}
% Distinguibilita e utilita delle feature per la navigazione, senza
% presentare DINOv2 come un classificatore esplicito dei materiali.

\section{Navigazione nello scenario con percorsi speculari}
\label{sec:esperimenti-speculari}
% Configurazione finale small specular map: scelta del percorso,
% successo e costo energetico effettivo.

\section{Navigazione nello scenario esteso}
\label{sec:esperimenti-esteso}
% Configurazione scelta per big general map. Distinguere nuove missioni
% su una mappa nota e generalizzazione a nuove mappe.

\section{Esperimenti sul robot reale}
\label{sec:esperimenti-reale}
% Sezione subordinata alle prove effettivamente completate:
% misura energetica, condizioni e limiti del confronto con la simulazione.

\section{Discussione complessiva dei risultati e dei casi di fallimento}
\label{sec:esperimenti-discussione}
% Integrare energia, affidabilita e durata delle missioni; analizzare
% i fallimenti sulla base dei dati. Risposta complessiva alla domanda
% di ricerca e limiti generali appartengono al capitolo 5.
